\begin{lemma}[Exit-coloring lemma]
Let $k\geq\Delta(G)+1$.  Let $T$ be a tree-sequence with first edge $e_0$,
and let $\varphi$ be a proper partial $k$-edge-coloring leaving exactly
$e_0$ uncolored.  Suppose that $T$ is closed with respect to $\varphi$ and
that $f\in\partial_G(V(T))$ has ends $u\notin V(T)$ and $v\in V(T)$ and
color $\delta=\varphi(f)$.  Suppose there is a proper partial coloring
$\pi$, again leaving exactly $e_0$ uncolored, which is
$(T,C\cup\{\delta\},\varphi)$-stable and for which some color is missing
both at $u$ and somewhere on $T$.  Then, for every
$\alpha\in\overline\varphi(v)$, there is a proper partial coloring
$\sigma$, leaving exactly $e_0$ uncolored and
$(T,C\cup\{\delta\},\varphi)$-stable, such that $v$ is a
$(T,\sigma,\{\alpha,\delta\})$-exit.
\end{lemma}
\begin{proof}
Choose the promised coloring $\pi$ and a color
$\beta\in\overline\pi(u)\cap\overline\pi(T)$.  By stability condition (ii)
in \noderef{StableColoring},
$\overline\pi(w)=\overline\varphi(w)$ at every $w\in V(T)$.  Hence
$\beta\in\overline\varphi(T)$.  The chosen
$\alpha\in\overline\varphi(v)$ also belongs to
$\overline\varphi(T)$.

Both $\alpha$ and $\beta$ are closed on $T$ under $\varphi$, because $T$
is closed.  They remain closed under $\pi$.  Indeed, if a boundary edge
$g$ had $\pi(g)$ equal to either color, then $g$ is incident with $T$ and
that color lies in $\overline\varphi(T)$.  The alternate stability
criterion \noderef{sc1} would give $\varphi(g)=\pi(g)$, contradicting
closure under $\varphi$ as defined in \noderef{ColoredVertexSetClosed}.

Define
\[
             \sigma=\pi/(G-T,\alpha,\beta)
\]
using \noderef{OutsideTreeColorSwap}.  Thus every edge incident with $T$ is
fixed, while $\alpha$ and $\beta$ are interchanged on all other edges.  We
check properness.  Two adjacent edges both outside $T$ undergo the same
transposition, and two adjacent edges both incident with $T$ are both fixed.
If exactly one is incident with $T$, their common vertex is outside $T$;
the incident edge is a boundary edge.  Such a boundary edge has neither
color $\alpha$ nor $\beta$ by the preceding paragraph, so the transposition
on the other edge cannot create a new equality.  Thus properness of $\pi$
implies properness of $\sigma$.  The only uncolored edge $e_0$ is an edge
of $T$ and is fixed, while the swap sends colors to colors, so $e_0$ remains
the unique uncolored edge.

Every edge incident with $T$ is fixed and consequently every missing-color
set at a vertex of $T$ is fixed.  Therefore $\sigma$ is
$(T,C\cup\{\delta\},\pi)$-stable directly from
\noderef{StableColoring}.  Since $\pi$ is
$(T,C\cup\{\delta\},\varphi)$-stable, transitivity in \noderef{sc2} shows
that $\sigma$ is $(T,C\cup\{\delta\},\varphi)$-stable.

The color $\alpha$ is missing at $u$ under $\sigma$.  To see this, every
$\alpha$- or $\beta$-colored edge incident with $u$ but not with $T$ is
swapped, and no boundary edge incident with $u$ has either color.  Hence an
$\alpha$-edge at $u$ under $\sigma$ would correspond to a $\beta$-edge at
$u$ under $\pi$, contradicting the choice
$\beta\in\overline\pi(u)$.  Also $\alpha$ remains missing at $v$, because
stability makes it missing there under $\pi$ and all edges incident with
$T$ are fixed.

Finally $f$ is incident with $T$, and its $\varphi$-color $\delta$ lies in
$C\cup\{\delta\}$.  Stability of $\pi$ and the outside swap therefore give
\[
                    \sigma(f)=\pi(f)=\varphi(f)=\delta.
\]
At each of $u$ and $v$, properness permits no second incident
$\delta$-edge, while $\alpha$ is missing at both vertices.  Thus in
$I_G(\sigma;\alpha,\delta)$ the three-node walk
\[
             u,\ f,\ v
\]
is a simple path, has incidence-graph length two, and cannot be extended at
either end.  Its only vertex in $V(T)$ is $v$, because $u\notin V(T)$, and
$\alpha$ is missing at its outside end $u$.  By
\noderef{TreeExitPath} and \noderef{TreeExit}, it is a $T$-exit path and
$v$ is the required exit.  The constructed $\sigma$ has all asserted
properties.
\end{proof}
